\documentclass[11pt,a4paper]{article}
\usepackage[margin=1in]{geometry}
\usepackage{amsmath,amssymb,amsthm}
\usepackage{algorithm}
\usepackage{algpseudocode}
\usepackage{booktabs}
\usepackage{hyperref}

\theoremstyle{definition}
\newtheorem{definition}{\textbf{Definition}}[section]
\newtheorem{theorem}{\textbf{Theorem}}[section]
\newtheorem{lemma}[theorem]{\textbf{Lemma}}
\newtheorem{remark}{\textbf{Remark}}[section]

\title{\textbf{Curriculum Learning and Pacing Functions in LLM Pretraining}}
\author{\textbf{Team Scaibu} \\ \small Email: \texttt{jaydeep.9664920749@gmail.com} \\ \small \textbf{Architecture and Core Engine Team}}
\date{\textbf{October 2026}}

\begin{document}
\maketitle

\section{Definitions and Cognitive Mental Models}

\subsection{Formal Mathematical Definition}
\textbf{Definition (Non-Stationary Data Scheduling and Pacing Functions):}
In classical empirical risk minimization, training batches are sampled uniformly and independently from an invariant data distribution $\mathcal{D}$. However, in complex high-dimensional sequence modeling, exposing randomly initialized models to extremely long, noisy, or linguistically intricate examples induces unstable gradient variance and slow optimization dynamics. \textbf{Curriculum Learning} imposes an explicit non-stationary data schedule ordered by difficulty, presenting simple, clean concepts early before progressively introducing harder, more nuanced distributions.

Let each training example $x \in \mathcal{D}$ possess a difficulty metric $d(x) \in \mathbb{R}_{\ge 0}$ (e.g., sequence length $|x|$, loss under a small reference model $\mathcal{L}_{\text{ref}}(x)$, or parse-tree depth). A \textbf{Pacing Function} $g: [0, T_{\text{pace}}] \to [\gamma_0, 1.0]$ governs the fraction of easiest training examples available to the optimizer at training step $t$:
{\boldmath
\begin{equation}
\mathcal{D}_t = \left\{ x \in \mathcal{D} \;\middle|\; d(x) \le F_d^{-1}(g(t)) \right\}
\end{equation}
}
where $F_d(\cdot)$ is the cumulative distribution function of sample difficulties, $\gamma_0 \in (0, 1]$ is the initial data fraction, and $g(t)$ is monotonically non-decreasing.

Common mathematical forms of pacing schedules include:
\begin{enumerate}
    \item \textbf{Linear Pacing}:
    {\boldmath
    \begin{equation}
    g_{\text{linear}}(t) = \min\left(1.0, \; \gamma_0 + (1 - \gamma_0)\frac{t}{T_{\text{pace}}}\right)
    \end{equation}
    }
    \item \textbf{Root Pacing} (rapid initial expansion):
    {\boldmath
    \begin{equation}
    g_{\text{root}}(t) = \min\left(1.0, \; \gamma_0 + (1 - \gamma_0)\sqrt{\frac{t}{T_{\text{pace}}}}\right)
    \end{equation}
    }
    \item \textbf{Geometric / Exponential Pacing}:
    {\boldmath
    \begin{equation}
    g_{\text{exp}}(t) = \min\left(1.0, \; \gamma_0 \cdot \left(\frac{1}{\gamma_0}\right)^{t / T_{\text{pace}}}\right)
    \end{equation}
    }
\end{enumerate}

\subsection{Expanded Non-Technical Definition (Intuition for Anyone)}
\textbf{Intuitive Conceptual Definition:}
\begin{enumerate}
    \item \textbf{How Humans Learn}: We do not hand a 5-year-old child a 600-page textbook on quantum electrodynamics on their first day of kindergarten. We start with the alphabet, advance to short sentences, then paragraphs, and only later tackle advanced calculus.
    \item \textbf{Avoiding Brain Overload}: When an artificial neural network starts with random weights, reading a 10,000-word convoluted legal contract produces chaotic gradients that knock its weights wildly out of alignment. Starting with simple, clear grammar helps it build foundational language concepts smoothly.
    \item \textbf{The Pacing Valve}: A pacing function acts like a controlled valve. At step zero, it lets through only the gentlest 20\% of the data. Over the next 50,000 steps, it opens wider and wider until the entire dataset is flowing.
    \item \textbf{Annealing Quality at the Finish Line}: In modern LLM training (Llama 3, DeepSeek), curriculum learning also runs in reverse at the very end (decay/annealing): during the final 10\% of training, the model is fed ultra-high-quality textbooks and math to polish its reasoning.
\end{enumerate}

\subsection{Child-Friendly Mental Model (Physical Metaphor)}
Imagine learning to swim in an indoor pool:
\begin{itemize}
    \item \textbf{The Disaster Method}: Tossing a toddler straight into a 10-meter deep whirlpool with 3-meter waves!
    \item \textbf{The Curriculum Method}:
    \begin{enumerate}
        \item Start in the 1-foot shallow wading pool with floaties (short, simple sentences).
        \item Move to the 4-foot lane to practice kicking and breathing (medium paragraphs).
        \item Finally, swim freely across the deep Olympic pool (massive documents and complex code).
    \end{enumerate}
    \item \textbf{Result}: The swimmer builds strong swimming muscles and confidence without ever swallowing water or panicking.
\end{itemize}

\section{Core Mathematical Equations and Definitions}

\subsection{Pacing Function Schedule}
{\boldmath
\begin{equation}
g(t) = \min\left(1.0, \; \gamma_0 + (1 - \gamma_0) \cdot \left(\frac{t}{T_{\text{pace}}}\right)^p\right)
\end{equation}
}
\begin{itemize}
    \item \textbf{Formal Definition}: Generalized polynomial pacing function mapping step $t$ to accessible percentile of training data.
    \item \textbf{What It Means}: Controls the velocity at which harder data examples are admitted into training batches ($p = 1$ is linear, $p = 0.5$ is root pacing).
    \item \textbf{Why It Is Important}: Prevents catastrophic early gradient spikes while guaranteeing full dataset coverage before warmup ends.
\end{itemize}

\subsection{Empirical Difficulty Quantile Threshold}
{\boldmath
\begin{equation}
d_{\text{thresh}}(t) = \inf \left\{ d \in \mathbb{R} \;\middle|\; \frac{1}{N}\sum_{i=1}^N \mathbb{I}(d(x_i) \le d) \ge g(t) \right\}
\end{equation}
}
\begin{itemize}
    \item \textbf{Formal Definition}: The empirical cutoff value filtering training samples based on measured difficulty scores.
    \item \textbf{What It Means}: Any sample with measured difficulty above this threshold is temporarily masked from optimization.
    \item \textbf{Why It Is Important}: Converts abstract pacing proportions into concrete sample-level eligibility filtering masks.
\end{itemize}

\section{Rigorous Mathematical Analysis Checklist}

\begin{itemize}
    \item \textbf{Notation}: $d(x)$ (difficulty score), $g(t)$ (pacing fraction), $\gamma_0$ (initial data fraction), $T_{\text{pace}}$ (pacing duration), $\mathcal{D}_t$ (active training subset).
    \item \textbf{Epistemic Status}: Introduced by Bengio et al. (ICML, 2009); universally adapted in LLM pretraining (progressive length scaling and data annealing).
    \item \textbf{Dependency Graph}: Dataset $\mathcal{D} \to \text{Difficulty Scoring } d(x) \to \text{Sort} \to g(t) \to \text{Active Batch Selection} \to \nabla_\theta \mathcal{L}$.
    \item \textbf{Dimensions}: Total dataset size $N$, active subset size $N_t = \lceil g(t) N \rceil$.
    \item \textbf{Regimes}: Language model pretraining, long-context extension, mathematical reasoning curriculum, speech recognition.
    \item \textbf{Invariants}: If $t \ge T_{\text{pace}}$, then $g(t) \equiv 1.0$, $d_{\text{thresh}} = \max_i d(x_i)$, and all $N$ samples are eligible.
    \item \textbf{Isomorphisms}: Curriculum learning is isomorphic to continuation methods in non-convex optimization and simulated annealing.
\end{itemize}

\section{Theoretical Theorems and Formal Proofs}

\begin{theorem}[Monotonic Convex Hull Expansion of Curriculum Datasets]
Let $g(t)$ be a monotonically non-decreasing pacing function: $g(t_1) \le g(t_2)$ for all $t_1 \le t_2$. Then the sequence of active datasets forms a monotonic filtration:
{\boldmath
\begin{equation}
\mathcal{D}_{t_1} \subseteq \mathcal{D}_{t_2} \subseteq \mathcal{D}
\end{equation}
}
and the convex hull of active data distributions expands monotonically: $\operatorname{conv}(\mathcal{D}_{t_1}) \subseteq \operatorname{conv}(\mathcal{D}_{t_2})$.
\end{theorem}

\begin{proof}
Let $x \in \mathcal{D}_{t_1}$. By definition of the active subset:
{\boldmath
\begin{equation}
d(x) \le F_d^{-1}(g(t_1))
\end{equation}
}
Because $F_d(\cdot)$ is a non-decreasing cumulative distribution function, its generalized inverse $F_d^{-1}(\cdot)$ is monotonically non-decreasing.
Since $t_1 \le t_2$ and $g$ is non-decreasing, $g(t_1) \le g(t_2)$, which implies:
{\boldmath
\begin{equation}
F_d^{-1}(g(t_1)) \le F_d^{-1}(g(t_2))
\end{equation}
}
Therefore:
{\boldmath
\begin{equation}
d(x) \le F_d^{-1}(g(t_1)) \le F_d^{-1}(g(t_2)) \implies x \in \mathcal{D}_{t_2}
\end{equation}
}
This proves $\mathcal{D}_{t_1} \subseteq \mathcal{D}_{t_2}$.
Because the convex hull operation is monotone with respect to set inclusion ($A \subseteq B \implies \operatorname{conv}(A) \subseteq \operatorname{conv}(B)$), it follows directly that $\operatorname{conv}(\mathcal{D}_{t_1}) \subseteq \operatorname{conv}(\mathcal{D}_{t_2})$.
\end{proof}

\begin{theorem}[Variance Reduction Bound Under Initial Low-Difficulty Warmup]
Let gradient noise covariance be proportional to sample loss: $\mathbb{E}[\|\nabla \ell(x) - \nabla L\|^2] \le C \cdot d(x)$. Then the initial gradient variance under curriculum learning is bounded by $\gamma_0 \cdot \operatorname{Var}_{\text{uniform}}$, strictly reducing initial optimization noise.
\end{theorem}

\begin{proof}
Under uniform sampling from $\mathcal{D}$:
{\boldmath
\begin{equation}
\operatorname{Var}_{\text{uniform}} \le C \cdot \mathbb{E}_{x \sim \mathcal{D}}[d(x)] = C \int_0^1 F_d^{-1}(u) \, du
\end{equation}
}
Under curriculum pacing at step $t=0$, samples are drawn uniformly from the easiest fraction $\gamma_0$:
{\boldmath
\begin{equation}
\operatorname{Var}_{\text{curr}}(0) \le C \cdot \mathbb{E}_{x \sim \mathcal{D}_0}[d(x)] = C \frac{1}{\gamma_0} \int_0^{\gamma_0} F_d^{-1}(u) \, du
\end{equation}
}
Because $F_d^{-1}(u)$ is a non-decreasing function:
{\boldmath
\begin{equation}
\frac{1}{\gamma_0} \int_0^{\gamma_0} F_d^{-1}(u) \, du \le F_d^{-1}(\gamma_0) \le \frac{1}{1 - \gamma_0} \int_{\gamma_0}^1 F_d^{-1}(u) \, du
\end{equation}
}
Consequently, $\mathbb{E}_{x \sim \mathcal{D}_0}[d(x)] < \mathbb{E}_{x \sim \mathcal{D}}[d(x)]$ whenever the dataset difficulty is non-constant.
Therefore, gradient variance during the critical initial optimization phase is strictly smaller than under uniform training, preventing catastrophic early divergence.
\end{proof}

\section{Foundational Architectural Questions and Epistemic Meta-Questions}

\subsection{Architectural Question 1: How to Measure Difficulty at Scale}
\textbf{Question:} How is sample difficulty practically computed for trillions of tokens in LLM pretraining?\\
\textbf{Answer:} Computing complex metrics online is too slow. Two efficient offline heuristics dominate: (1) Sequence length, training first on 2K tokens, then 4K, 8K, and finally 32K/128K; and (2) Cross-entropy loss under a small fast reference model (e.g., a 1B parameter proxy model evaluated during dataset tokenization).

\textbf{Meta-Question:} Can curriculum learning ever harm final model accuracy?\\
\textbf{Answer:} Yes. If a model overfits to simple examples during a prolonged curriculum warmup, it can suffer from plasticity loss or catastrophic forgetting of complex reasoning. Pacing schedules must not dwell excessively on easy data, and the final 50--80\% of training should operate over the full unconstrained distribution.

\subsection{Architectural Question 2: Standard Curriculum vs Anti-Curriculum}
\textbf{Question:} When does an anti-curriculum (hard examples first) outperform easy-first training?\\
\textbf{Answer:} In fine-tuning tasks where the model already has strong pre-existing representations, presenting difficult, high-information examples first forces rapid weight adaptation, whereas simple repetitive examples cause gradient wasting. Easy-first is optimal for pretraining from scratch; hard-first or mixed curriculum is often superior for domain adaptation.

\textbf{Meta-Question:} How is curriculum learning implemented in mixture-of-experts (MoE) routing?\\
\textbf{Answer:} In MoE training, routing entropy is paced: router auxiliary balance loss is kept high early in training to force balanced expert exploration, then gently relaxed later in training to allow specialized expert differentiation.

\section{Algorithmic Implementation Specification}

\begin{algorithm}[H]
\caption{Curriculum Pacing and Batch Eligibility Filter}
\begin{algorithmic}[1]
\Require Sample difficulties $\mathbf{d} \in \mathbb{R}^N$, step $t$, pacing duration $T_{\text{pace}}$, mode $pacing\_type$, initial fraction $\gamma_0$
\Ensure Eligible indices $\mathcal{I}_{\text{elig}}$, threshold $d_{\text{thresh}}$, active fraction $f$, sorted indices $\mathcal{I}_{\text{sort}}$
\State $progress \gets \min(1.0, t / T_{\text{pace}})$
\If{$progress \ge 1.0$}
    \State $f \gets 1.0$
\ElsIf{$pacing\_type == \text{"linear"}$}
    \State $f \gets \gamma_0 + (1.0 - \gamma_0) \cdot progress$
\ElsIf{$pacing\_type == \text{"root"}$}
    \State $f \gets \gamma_0 + (1.0 - \gamma_0) \cdot \sqrt{progress}$
\ElsIf{$pacing\_type == \text{"exponential"}$}
    \State $f \gets \gamma_0 \cdot (1.0 / \gamma_0)^{progress}$
\Else
    \State $f \gets \gamma_0 + (1.0 - \gamma_0) \cdot progress$
\EndIf
\State $f \gets \max(\gamma_0, \min(1.0, f))$
\State Sort pairs $(i, \mathbf{d}[i])$ ascending by difficulty to produce sorted indices $\mathcal{I}_{\text{sort}}$
\State $count \gets \max(1, \min(N, \lceil f \cdot N \rceil))$
\State $\mathcal{I}_{\text{elig}} \gets \mathcal{I}_{\text{sort}}[0 : count]$
\State $d_{\text{thresh}} \gets \mathbf{d}[\mathcal{I}_{\text{elig}}[count - 1]]$
\State \Return $(\mathcal{I}_{\text{elig}}, d_{\text{thresh}}, count / N, \mathcal{I}_{\text{sort}})$
\end{algorithmic}
\end{algorithm}

\section{Big-$\Theta$ Complexity Analysis}

\begin{itemize}
    \item \textbf{Sorting Time Complexity}: $\Theta(N \log N)$ to rank sample difficulties.
    \item \textbf{Pacing Fraction Evaluation}: $\Theta(1)$ scalar math operations.
    \item \textbf{Subsetting and Slicing Time}: $\Theta(\lceil f N \rceil)$ index allocation.
    \item \textbf{Space Complexity}: $\Theta(N)$ buffer for sorted index arrays.
    \item \textbf{Parallel Sorting Depth}: $\mathcal{D} = \Theta(\log^2 N)$ under parallel bitonic sorting networks.
\end{itemize}

\section{Single Canonical Repository Reference}

The complete deterministic scalar implementation, verification suite, and unit test benchmarks are published at:
\begin{center}
\url{https://github.com/Chief-Strategist-J/llm-obs-infra}
\end{center}

\end{document}
